%%%%%%%%%%%%%%%%%%%%%%%%%%%%%%%%%%%%%%%%%%%%%%%%%%%%%%%%%%%%%%%%%%%%%%%%%%%
% Handout und Uebungsblatt zu Lehreinheit 12
% Im Skript am Kapitelende, in handouts.tex als Druckvorlage.
%%%%%%%%%%%%%%%%%%%%%%%%%%%%%%%%%%%%%%%%%%%%%%%%%%%%%%%%%%%%%%%%%%%%%%%%%%%
\begin{handout}{Präsentation und Glaubwürdigkeit}

\textbf{Worum es heute ging.} Zwei Vortragsausschnitte: Der eine wirkte
überzeugender, der andere war fachlich belastbarer. Was überzeugt, ist
nicht dasselbe wie das, was stimmt. Das gilt für fremde Vorträge -- und
für Ihre eigenen.

\textbf{Kernbegriffe}
\begin{description}
  \item[Credibility Cue] Signal, an dem ein Publikum Glaubwürdigkeit
    abliest (Stimme, Sicherheit, Titel, Fachwörter, gezeigte Daten).
  \item[CREDS] Credibility Enhancing Displays: Handlungen mit Kosten oder
    Risiko, die Überzeugung beweisen sollen \parencite{henrich2009}. Sie
    zeigen, wie überzeugt jemand ist, nicht, ob er recht hat.
  \item[Dr.-Fox-Effekt] Ein mitreißender Vortrag ohne Inhalt wird gut
    bewertet; Charisma hebt die gefühlte, nicht die tatsächliche
    Lernleistung \parencite{naftulin1973drfox,peerbabad2014drfox}.
  \item[Arbeitsgedächtnis] hält etwa 3--5 Einheiten zugleich
    \parencite{cowan2001magical}; Folien sollen es nicht unnötig belasten.
  \item[Dual Coding] Wort und Bild werden getrennt verarbeitet und
    gemeinsam besser behalten \parencite{paivio1986mental}.
\end{description}

\textbf{Der 5-Minuten-Vortrag}

\begin{center}\small
\begin{tabular}{p{0.18\textwidth}p{0.13\textwidth}p{0.55\textwidth}}
\toprule
\textbf{Teil} & \textbf{Zeit} & \textbf{Leitfrage} \\
\midrule
Frage, Lücke & ca.\ 1:45 & Was will ich wissen, warum ist das offen? \\
Methode & ca.\ 1:00 & Wie habe ich es geprüft? \\
Befunde & ca.\ 1:15 & Was kam heraus? (dann: Pause) \\
Grenzen & ca.\ 0:30 & Was kann ich (noch) nicht sagen? \\
Botschaft & ca.\ 0:30 & Ein Satz, der bleibt. \\
\bottomrule
\end{tabular}
\end{center}

\textbf{Folienregeln.} Überschrift als Aussage, darunter der Beleg
\parencite{alley2013craft}. Alles weg, was nicht zur Botschaft gehört;
nicht vorlesen, was auf der Folie steht \parencite{mayer2009multimedia}.
Beschriftung an die Daten, nicht in eine ferne Legende. Achsen nicht
abschneiden, ohne es zu sagen.

\begin{merkbox}
Glaubwürdigkeit entsteht im Publikum, Belastbarkeit liegt in der Aussage.
Wer Grenzen sachlich benennt, macht sich angreifbar -- und gerade dadurch
glaubwürdig.
\end{merkbox}

\textbf{Zum Weiterlesen:} \textcite{alley2013craft}, Teil zu Folien;
\textcite{henrich2009}.
\end{handout}

\begin{uebungsblatt}{Überzeugen und belegen}

\aufgabe{Zwei Vorträge im Vergleich}[in der Stunde, Einzelarbeit, 5 Min.]
Bewerten Sie jeden Ausschnitt direkt nach dem Ansehen (1 = gar nicht,
5 = sehr) und nennen Sie je einen Grund.

\begin{center}
\renewcommand{\arraystretch}{2.0}
\begin{tabular}{|p{0.12\textwidth}|p{0.1\textwidth}|p{0.25\textwidth}|p{0.1\textwidth}|p{0.25\textwidth}|}
\hline
 & \textbf{überzeugend} & \textbf{weil \ldots} & \textbf{belastbar} & \textbf{weil \ldots} \\
\hline
Ausschnitt A & & & & \\ \hline
Ausschnitt B & & & & \\ \hline
\end{tabular}
\end{center}

\aufgabe{Feedbackbogen zur 5-Minuten-Präsentation}[in der Stunde, während der Vorträge]
Halten Sie zu jedem Vortrag je einen Punkt fest. Beschreiben Sie konkret
(\enquote{Die Folie mit dem Balkendiagramm \ldots}), statt allgemein zu
loben oder zu tadeln. Den Bogen bekommen die Vortragenden.

\begin{center}\small
\renewcommand{\arraystretch}{3.2}
\begin{tabular}{|p{0.16\textwidth}|p{0.24\textwidth}|p{0.24\textwidth}|p{0.24\textwidth}|}
\hline
\textbf{Vortrag von} & \textbf{Was war klar?} & \textbf{Was hat überzeugt -- und war es belegt?} & \textbf{Was blieb offen?} \\
\hline
 & & & \\ \hline
 & & & \\ \hline
 & & & \\ \hline
 & & & \\ \hline
 & & & \\ \hline
\end{tabular}
\end{center}

\end{uebungsblatt}

\begin{uebungsblattb}{3}{Bewerten und Gestalten}

\aufgabe{Bewertungsraster der Präsentation}[zur Information; die Lehrperson bewertet danach]
Die Präsentation zählt 15\,\% der Modulnote. Inhaltliche Kriterien zählen
doppelt -- das ist die Konsequenz aus dem Dr.-Fox-Effekt: Ein mitreißender
Vortrag wird auch ohne Substanz gut bewertet.

\begin{center}\small
\renewcommand{\arraystretch}{1.5}
\begin{tabular}{|p{0.5\textwidth}|c|c|c|c|c|}
\hline
\textbf{Kriterium} & \textbf{0} & \textbf{1} & \textbf{2} & \textbf{3} & \textbf{Faktor} \\
\hline
Forschungsfrage klar und begründet & & & & & $\times 2$ \\ \hline
Methode nachvollziehbar dargestellt & & & & & $\times 2$ \\ \hline
Befunde und Grenzen transparent; Sicherheit passt zu den Daten & & & & & $\times 2$ \\ \hline
Aufbau und roter Faden; Zeit eingehalten ($\leq$ 5:00) & & & & & $\times 1$ \\ \hline
Folien: lesbar, eine Botschaft je Folie, Belege statt Dekoration & & & & & $\times 1$ \\ \hline
Vortrag und Umgang mit Fragen & & & & & $\times 1$ \\ \hline
\multicolumn{5}{|r|}{\textbf{Summe} (höchstens 27 Punkte)} & \\ \hline
\end{tabular}
\end{center}
{\small 0 = fehlt, 1 = in Ansätzen, 2 = überwiegend erfüllt, 3 = voll erfüllt.}

\aufgabe{Folienkritik}[in der Stunde, Kleingruppe, 15 Min.]
Die drei Folien stammen aus einem \emph{fiktiven} Vortrag über ein
Seminarexperiment; alle Daten sind erfunden. Analysieren Sie jede Folie
mit dem Raster und skizzieren Sie eine bessere Fassung.

\medskip
\noindent\begin{tikzpicture}[font=\sffamily]
  % Folie 1: Textwueste
  \draw[rwugray,thick] (0,0) rectangle (7.6,4.6);
  \node[anchor=north west,font=\sffamily\bfseries\small] at (0.15,4.5) {Folie 1: Methodik};
  \node[anchor=north west,text width=7.2cm,font=\sffamily\tiny,align=left] at (0.15,4.0) {%
    \textbullet\ Es wurde ein Experiment mit Studierenden der Hochschule durchgeführt, die freiwillig teilgenommen haben.\\
    \textbullet\ Die Teilnehmenden wurden in zwei Gruppen eingeteilt, wobei die Einteilung nach dem Zufallsprinzip erfolgte.\\
    \textbullet\ Gruppe A sah vor der Lerneinheit einen Zaubertrick, Gruppe B sah stattdessen ein kurzes Video ohne Überraschung.\\
    \textbullet\ Danach wurde in beiden Gruppen dieselbe Lerneinheit von 15 Minuten zum Thema Wahrscheinlichkeit gehalten.\\
    \textbullet\ Im Anschluss wurde ein Test mit 10 Fragen geschrieben, der von einer Person ausgewertet wurde.\\
    \textbullet\ Die auswertende Person wusste nicht, wer zu welcher Gruppe gehörte (verblindet).\\
    \textbullet\ Außerdem wurden Alter, Geschlecht, Studiengang und Vorwissen erhoben.\\
    \textbullet\ Eine Woche später wurde der Test wiederholt, allerdings nahmen nicht mehr alle teil.};
  % Folie 2: Diagramm mit abgeschnittener Achse
  \begin{scope}[xshift=8.2cm]
  \draw[rwugray,thick] (0,0) rectangle (7.6,4.6);
  \node[anchor=north west,font=\sffamily\bfseries\small] at (0.15,4.5) {Folie 2: Ergebnisse};
  \foreach \y/\l in {0.6/6{,}0,1.1/6{,}2,1.6/6{,}4,2.1/6{,}6,2.6/6{,}8,3.1/7{,}0,3.6/7{,}2}{
    \draw[black!25] (0.9,\y) -- (5.2,\y);
    \node[left,font=\sffamily\tiny] at (0.9,\y) {\l};}
  \fill[red!70] (1.3,0.6) rectangle (2.3,3.35);
  \fill[green!60!black] (2.9,0.6) rectangle (3.9,2.6);
  \draw[black,thick] (0.9,0.6) -- (5.2,0.6);
  \draw[black,thick] (0.9,0.6) -- (0.9,3.9);
  \node[font=\sffamily\tiny,anchor=west] at (5.4,3.4) {\textcolor{red!70}{$\blacksquare$} Serie 1};
  \node[font=\sffamily\tiny,anchor=west] at (5.4,3.1) {\textcolor{green!60!black}{$\blacksquare$} Serie 2};
  \node[font=\sffamily\tiny,anchor=west] at (5.4,2.8) {\textcolor{blue}{$\blacksquare$} Serie 3};
  \node[font=\sffamily\tiny] at (3,0.3) {Mittelwert Punkte};
  \end{scope}
\end{tikzpicture}

\medskip
\noindent\begin{tikzpicture}[font=\sffamily]
  % Folie 3: Botschaft passt nicht zu den Daten
  \draw[rwugray,thick] (0,0) rectangle (7.6,4.6);
  \node[anchor=north west,font=\sffamily\bfseries\small,text width=7.2cm] at (0.15,4.5)
    {Folie 3: Staunen verbessert das Lernen deutlich!};
  \fill[rwuviolet!50] (1.4,0.6) rectangle (2.4,2.64);
  \fill[rwucyan] (3.2,0.6) rectangle (4.2,2.73);
  \draw[thick] (1.9,2.19) -- (1.9,3.09); \draw[thick] (1.75,2.19) -- (2.05,2.19); \draw[thick] (1.75,3.09) -- (2.05,3.09);
  \draw[thick] (3.7,2.28) -- (3.7,3.18); \draw[thick] (3.55,2.28) -- (3.85,2.28); \draw[thick] (3.55,3.18) -- (3.85,3.18);
  \draw (1,0.6) -- (4.6,0.6);
  \node[font=\sffamily\tiny] at (1.9,0.35) {ohne Trick: 6,8};
  \node[font=\sffamily\tiny] at (3.7,0.35) {mit Trick: 7,1};
  \node[font=\sffamily\tiny,align=left,anchor=west] at (4.8,2.2) {je $n = 12$\\$p = 0{,}41$};
  \node[draw,dashed,font=\sffamily\tiny,align=center,minimum width=1.6cm,minimum height=1cm] at (6.3,3.4) {[Clipart:\\Glühbirne]};
\end{tikzpicture}\hfill
\begin{minipage}[b]{0.45\textwidth}\small
Folie~2 zeigt dieselben Mittelwerte wie Folie~3: roter Balken 7,1
(mit Trick), grüner Balken 6,8 (ohne Trick), jeweils von 10 möglichen
Punkten.
\end{minipage}

\begin{center}\small
\renewcommand{\arraystretch}{2.4}
\begin{tabular}{|c|p{0.25\textwidth}|p{0.25\textwidth}|p{0.25\textwidth}|}
\hline
\textbf{Folie} & \textbf{Klarheit: Was ist die Botschaft?} & \textbf{Arbeitsgedächtnis: Was belastet unnötig?} & \textbf{Konsistenz: Passen Titel, Bild und Daten?} \\
\hline
1 & & & \\ \hline
2 & & & \\ \hline
3 & & & \\ \hline
\end{tabular}
\end{center}
Skizze einer besseren Fassung (eine Folie Ihrer Wahl):
\schreiblinien{4}

\aufgabe{Tagebuch des Staunens}[Hausaufgabe]
Wann haben Sie zuletzt jemandem geglaubt, weil er überzeugend
\emph{wirkte}? Welcher Credibility Cue war es (welches Signal, an dem
Sie Glaubwürdigkeit abgelesen haben), und was hätte ihn als Beleg
gedeckt?

\begin{lehrhinweis}[Erwartungshorizont]
\textbf{Aufgabe~1:} Erwartet wird, dass mindestens ein Teil der Gruppe A
auf beiden Skalen höher bewertet; die Begründungen für
\enquote{belastbar} nennen dann oft Merkmale der Performance. Genau das ist
der Diskussionsanlass.

\textbf{Aufgabe~2:} Gute Rückmeldungen sind konkret und trennen bei
\enquote{überzeugt} zwischen Wirkung und Beleg.

\textbf{Aufgabe~4, Folie~1:} Keine Botschaft (Titel ist ein Thema),
acht ganze Sätze überfordern das Arbeitsgedächtnis, und die Vortragende
wird sie vorlesen (Redundanz). Besser: Titel als Aussage (\enquote{Zwei
zufällig gebildete Gruppen, verblindete Auswertung}), ein Ablaufschema mit
drei Kästchen, Details mündlich. Der wichtige Hinweis auf den Schwund bei
der Wiederholung gehört zu den Limitationen. \textbf{Folie~2:} Die
Achse beginnt bei 6,0 und lässt einen Unterschied von 0,3 Punkten riesig
erscheinen; Rot und Grün tragen Bewertungen
(\enquote{schlecht}, \enquote{gut}), die nicht gemeint sind, und sind für
Menschen mit Rot-Grün-Schwäche kaum zu unterscheiden; der Titel sagt
nichts, die Gruppen sind nicht beschriftet, es fehlt eine Achsenbeschriftung, und die Legende
nennt eine \enquote{Serie~3}, die es nicht gibt. Besser: Achse ab~0 oder
den Abschnitt deutlich kennzeichnen, Gruppen direkt beschriften, Streuung
zeigen. \textbf{Folie~3:} Der Titel behauptet einen deutlichen Effekt, die
Daten zeigen einen kleinen Unterschied mit sich überlappenden
Fehlerbalken bei $n = 12$ je Gruppe und $p = 0{,}41$; die Clipart ist
Dekoration. Besser: \enquote{Kein klarer Unterschied zwischen den Gruppen
($n = 12$ je Gruppe)} -- und als Grenze nennen, dass die Stichprobe für
kleine Effekte zu klein ist (LE~8).
\end{lehrhinweis}
\end{uebungsblattb}
